\chapter{Example programs}\label{ch:examples}

The \texttt{examples/} directory contains eleven programs, each commented line by line. They all run on
the bare machine, in supervisor mode, share the console helpers in \texttt{examples/lib/print.s}, and are
cross-checked between the emulator and the RTL by \texttt{make test}. This chapter explains what each one
demonstrates and shows its central code.

Nine of them also run unmodified under StrataOS (Chapter~\ref{ch:os}), in user mode, as bare-metal
programs: the disk carries each one assembled for I-LS \hex{8000}, and \texttt{-n N prog} runs it as $N$
time-shared threads on $N$ contiguous worker cores, which it sees as cores $0 \ldots N-1$ through
\csr{corebase} and \csr{corecount}; their \insn{ctake}, \insn{cpoll} and inbox \insn{tany} are emulated
by the dispatcher. \texttt{clock} (Section~\ref{sec:ex-clock}) and \texttt{gfx}
(Section~\ref{sec:ex-gfx}) are bare-metal only: they cut their device capabilities from the device
root \reg{c3}, which under StrataOS only the \texttt{dev} server holds.

\section{Conventions shared by the examples}
\begin{itemize}
\item Every core starts at \texttt{\_start}; cores not needed by an example \insn{halt} at once. The
  stack is where reset put \reg{sp}, at the top of the D-LS; no example assumes a Local Store size.
  What each needs is its image and its stack: every example runs with 32~KiB Local Stores;
  \texttt{clock}, \texttt{faults}, \texttt{hello}, \texttt{mandel} and \texttt{race} with 4~KiB of D-LS,
  \texttt{pipeline} with 8~KiB; all but \texttt{overlay} (its slots are at I-LS \hex{7700}) with a 4~KiB
  I-LS; \texttt{gfx}, \texttt{parsum}, \texttt{stream\_sum} and \texttt{swcache} need 32~KiB of D-LS for
  their buffers (\texttt{make test-ls}).
\item Output goes through \reg{c2} with \insn{tput} on tag \reg{tg15} (\texttt{print\_str},
  \texttt{print\_dec}, \texttt{print\_hex}, \texttt{print\_nl}, \texttt{print\_kv}). The helpers do not
  wait: reusing their buffer stalls until the previous \insn{tput} has read it.
\item Lines starting with \texttt{@} carry cycle counts or times; they differ between the emulator and the RTL and
  are ignored by \texttt{compare.sh}.
\end{itemize}

\section{\texttt{hello}: Local Stores, output, and a token}\label{sec:ex-hello}
Each core prints one line, in core order. Core $k$ waits with \insn{ctake c3, t0} for a token from core
$k-1$ only, prints, and passes a copy of its console capability on with \insn{cgive}. Because
\insn{cgive} waits for the giver's memory transfers, each line is completely out before the next core starts,
so the output is deterministic on any timing. The program is shown in Chapter~\ref{ch:intro}.

\begin{shell}
Hello from Strata core 0 of 4
Hello from Strata core 1 of 4
...
\end{shell}

\section{\texttt{stream\_sum}: streaming and double buffering}\label{sec:ex-stream}
Sums a 128~KiB array in T3 twice, in 8~KiB chunks. The plain version issues a \insn{tget} and immediately sums the buffer;
each load stalls until its bytes arrive. The double-buffered version prefetches chunk $k+1$ into one
buffer while summing chunk $k$ in the other. Both are correct by the as-if-synchronous rule; only the
cycle counts differ.

\begin{asm}
dbl_loop:
        addi  s4, s4, -1                ; this iteration consumes one chunk
        beqz  s4, dbl_sum               ; nothing left to prefetch
        tget  tg1, s8, c0, s2, s6       ; prefetch the next chunk into the other buffer
        add   s2, s2, s6
dbl_sum:
        mv    a0, s7                    ; sum the current chunk while the prefetch runs
        call  sum_chunk
        mv    t0, s7                    ; swap the two buffers
        mv    s7, s8
        mv    s8, t0
        bnez  s4, dbl_loop
\end{asm}

\section{\texttt{swcache}: a software-managed cache}\label{sec:ex-swcache}
Builds a 4-way, 32-set cache of 64-byte lines in the D-LS, describes it with \csr{ttag}, \csr{tdata} and
\csr{tgeom}, and reads 4000 pseudo-random entries of a 64~Ki-entry table in T3 through it, checking every
value. The lookup routine is shown in Section~\ref{sec:swcache}. Hit and miss counts are deterministic,
because the replacement policy is software-visible state, not timing.

\section{\texttt{overlay}: code loaded on demand}\label{sec:ex-overlay}
Three routines live in T3 (\texttt{.far}) and are called with \insn{xcall} through two 64-byte I-LS
slots. On a miss the handler loads the routine with \insn{tgetx}, maps it with \insn{xmap} and returns
with \insn{eret} to retry the call. It saves its temporaries to a D-LS area addressed from \reg{x0}
because a handler has no free register. It uses exact mapping (\csr{xslot} = 0); the StrataOS
\texttt{units} program shows slot mapping, far links and demand paging (Section~\ref{sec:os-paging}).
Under StrataOS the handler is the program's own user-mode handler: its \insn{csrw tvec} sets
\csr{utvec}, and its misses never reach the kernel.

\begin{asm}
handler:
        sd    t0, hsave+0(zero)         ; save temporaries without a base register
        ...
        csrr  t0, cause
        li    t1, 8                     ; 8 = xmiss
        bne   t0, t1, fatal
        csrr  t0, tval                  ; far address of the code
        cbase t1, c0
        sub   t0, t0, t1                ; its offset within c0
        ...
        tgetx tg2, t2, c0, t0, t3       ; copy the code into slot 0
        xmap  0, t2, c0, t0             ; map far address -> slot 0
        ...                             ; restore temporaries
        eret                            ; re-execute the xcall, which now hits
\end{asm}

\section{\texttt{faults}: the protection model, one trap at a time}\label{sec:ex-faults}
Installs a handler that records the cause and skips the trapping instruction, then commits ten
violations in turn: reading past a segment (\cause{bounds}), a misaligned transfer (\cause{xalign}),
using a capability after moving it (\cause{capinv}), writing through a share (\cause{perm}), a load
outside the D-LS (\cause{lsfault}), an atomic on ordinary memory (\cause{perm}), an undefined instruction
(\cause{illegal}), splitting beyond the end (\cause{capop}), overwriting a live capability
(\cause{capop}) and giving to a halted core (\cause{capop}).

\section{\texttt{pipeline}: zero-copy producer and consumer}\label{sec:ex-pipeline}
Core 0 fills 4~KiB buffers in T3 and gives each to core 1, which sums it and gives it back. Two buffers
circulate, so filling block $k+1$ overlaps summing block $k$. The data never moves between cores;
only capabilities do, and each \insn{cgive} both transfers ownership and guarantees the data is in
memory.

\begin{asm}
        tput  tg0, s5, c4, zero, s6     ; copy the block into the buffer c4 names
        li    t0, 1
        cgive c4, t0                    ; hand it over (waits for the tput; c4 becomes empty)
        ...
consumer:
        ctake c4                        ; wait for the next block's capability
        tget  tg0, s5, c4, zero, s6     ; bring it into our own Local Store
        li    t0, 0
        cgive c4, t0                    ; return the buffer at once (waits for the tget)
\end{asm}

\section{\texttt{parsum}: parallel reduction}\label{sec:ex-parsum}
Core 0 makes an 8-byte accumulator atomic with \insn{catomic}, gives every core a copy, splits the array
capability into one piece per core and gives each piece away. Every core sums its piece and adds the
partial sum with \insn{tamo.add}. Core 0 takes the pieces back \emph{in core order} with
\insn{ctake c5, s3} and \insn{cjoin}s them, so it ends owning the whole array again, regardless of the
order in which the cores finished.

\begin{asm}
par_take:
        bge   s3, s7, par_done          ; all pieces are back
        ctake c5, s3                    ; take a piece back from core s3 only
        cjoin c4, c5                    ; append it: c4 grows back to the whole array
        addi  s3, s3, 1
        j     par_take
par_done:
        tamo.add s8, c6, zero, zero     ; atomic add of 0 = atomic read of the total
\end{asm}

\section{\texttt{mandel}: one frame, many owners}\label{sec:ex-mandel}
The frame buffer in T3 is split into one band of rows per core. Each core renders into its own band
(a stray write into another band would trap) and gives the band back; core 0 rejoins the bands in core
order and prints the whole frame with one transfer. Arithmetic is fixed point with 24 fractional bits.
Each row is rendered into one D-LS buffer and sent with \insn{tput}; the next row's stores into that
buffer simply wait for the previous \insn{tput}.

\section{\texttt{race}: the races that cannot be written}\label{sec:ex-race}
Attempts the three classic races. A lost update is impossible because the shared counter is an atomic
segment, reachable only through \insn{tamo}; a plain \insn{tput} to it traps with \cause{perm}, and two
cores adding 1000 each always leave 2000. Use-after-give is impossible because giving empties the
register; the attempt traps with \cause{capinv}. A stale read is impossible because a \insn{cgive}
completes only after the giver's writes have landed.

\begin{shell}
core 0 writes a buffer it gave away   -> capinv
shared counter after 2 x 1000 adds    = 2000
core 1 plain store to the counter     -> perm
message from core 1: core 1 says hi through owned T3
\end{shell}

\section{\texttt{clock}: the timer, and events without interrupts}\label{sec:ex-clock}
Core 0 finds the timer in the platform table (entry type 6), cuts a capability for exactly its 32
bytes of device space out of the device root, and reads FREQ and NOW with 8-byte \insn{tget}s
(Section~\ref{sec:timer}). It then gives core 1 an 8-byte mailbox of system memory; core 1 sums the
squares of 1 to 300, writes the result into the mailbox and gives it back. Meanwhile core 0 waits for
four ticks 5~ms apart. Each tick is armed by writing DEADLINE and starting a read of WAIT, which stays
in flight until the timer reaches the deadline:

\begin{asm}
arm:
        add   s10, s5, t0               ; s10 = absolute deadline: no drift from late wakeups
        la    t2, dlbuf
        sd    s10, 0(t2)
        li    t3, R_DL                  ; DEADLINE
        li    t4, 8                     ; 8 bytes; the value takes effect when byte 7 lands
        tput  tg2, t2, c5, t3, t4       ; DEADLINE = s10
        la    t2, wbuf                  ; where NOW goes when the wait is over
        li    t3, R_WAIT
        tget  tg1, t2, c5, t3, t4       ; start waiting: tg1 stays busy until NOW >= DEADLINE
        ret                             ; return at once
\end{asm}

The event loop waits for whichever comes first, the WAIT completing (tag 1 idle) or a capability
arriving from core 1 (bit 16), with one \insn{tany}:

\begin{asm}
events:
        tany  t0, s7                    ; s7 = bit 1 (tag 1) | bit 16 (a gift)
        srli  t1, t0, 16
        andi  t1, t1, 1
        beqz  t1, no_gift
        li    t0, 1
        ctake c7, t0                    ; the mailbox, from core 1 only
        ...
no_gift:
        andi  t1, s8, 2                 ; tag 1: the WAIT completed, NOW is in wbuf
\end{asm}

There is no interrupt and no handler: a device event is a transfer completing, and waiting for several
events is waiting for several tags. Depending on timing, core 1's line comes before or after the first
tick; the program checks that every WAIT ended at or after its deadline. Times are printed on
\texttt{@} lines.

\begin{shell}
timer at device offset 12288, 1000000 ticks per second
FREQ register says 1000000
core 1 finished while we waited: 1^2 + ... + 300^2 = 9045050
tick 1
...
every WAIT ended at or after its DEADLINE: yes
\end{shell}

\section{\texttt{gfx}: a picture in the framebuffer, one band per core}\label{sec:ex-gfx}
Renders the Mandelbrot set in colour into the framebuffer, on every core. Nothing about the display is
built in: core 0 cuts a window onto the platform table out of the device root, removes its write
permission with \insn{crestrict}, and gives each core a copy, so every core reads the framebuffer entry
(type 5) itself: the base address in system memory, and width, height, stride and pixel format from
its info word (Section~\ref{sec:fbctl}). The program draws whatever size it is told, for example with
\texttt{strata-emu -r 640x480}.

Core 0 cuts the visible part of the framebuffer, height $\times$ stride bytes, out of its memory
capability and deals it out in bands of rows:

\begin{asm}
        csplit c6, c0, s2               ; c6 = system memory from the framebuffer on
        csplit c7, c6, s3               ; c6 = exactly the visible framebuffer
        csplit c5, c6, s9               ; c6 = band 0, c5 = the rest
give:
        ...
        csplit c8, c5, t2               ; c5 = core s8's band, c8 = the rest
        cgive c5, s8                    ; core s8 owns it now
        cmov  c5, c8
\end{asm}

Each core computes rows of points at half resolution in fixed point (24 fraction bits), stores two
pixels per point as one doubleword into one of two D-LS row buffers, and sends the row twice with
\insn{tput}, a 2\,x\,2 block per point, while the next row is computed into the other buffer. A core can
write only its own band. When it is done it gives the band back; \insn{cgive} completes only after the
core's transfers into it have landed, so once core 0 has taken every band back, in core order, and
rejoined them with \insn{cjoin}, the picture in memory is complete. The display scans memory out
live, so the bands appear as they are drawn.

Core 0 then cuts a capability for the framebuffer controller's two ports and writes PRESENT, ``this
frame is complete''; the simulators capture this frame for \texttt{-F}. Waiting for the display is
waiting for a transfer: a read of VSYNC completes at the next 60~Hz refresh with the frame number.

\begin{asm}
        tput  tg0, s1, c9, zero, t1     ; PRESENT (the value written is ignored)
        ...
vsync:
        li    t0, R_VSYNC
        li    t1, 8
        tget  tg1, s1, c9, t0, t1       ; completes at the next refresh
        ...
        ld    a0, 0(s1)                 ; the frame number: this load waits for the refresh
\end{asm}

\begin{shell}
framebuffer: 1280 x 800, 5120 bytes per row, format 1, at 0x0000000003000000
rendered by 8 cores; presented, waited for 3 refreshes
\end{shell}

\texttt{tools/test-fb.sh} runs \texttt{gfx} on the emulator and the RTL and requires the two presented
frames to be byte-identical; \texttt{make show-gfx} keeps its window open after it halts.
